\documentclass[a4paper,11pt]{article}
\usepackage{nsi-preamble}
\newcommand{\cor}[4]{\subsection*{Exercice #1 (#2)}\begin{itemize}[nosep]\item \textbf{Réponse} : #3\item \textbf{Cas limite} : #4\end{itemize}}
\begin{document}
\nsiheader{Corrigé — Dichotomie, glouton et k-NN}{Première}{P13}{1}{P-ALGO-03, P-ALGO-04, P-ALGO-05}{BO 2019}{needs\_review}

\section*{Corrigé du TD}
\cor{1}{P-ALGO-04}{milieux $18$ puis $37$ $\rightarrow$ indice $4$.}{cible absente.}
\cor{2}{P-ALGO-04}{$V = \text{droite} - \text{gauche} + 1$ décroît de $6$ à $3$ $\rightarrow$ terminaison.}{\texttt{cible=38} $\rightarrow$ $V$ décroît $6\rightarrow 3\rightarrow 1\rightarrow 0$, arrêt sans trouver.}
\cor{3}{P-ALGO-05}{$28 = 10+10+5+2+1$ (5 pièces).}{pièce 1 absente $\rightarrow$ le glouton peut échouer (montant~$3$ avec \texttt{[5,2]}).}
\cor{4}{P-ALGO-03}{$k=3$ $\rightarrow$ rouge~2, bleu~1 $\rightarrow$ classe rouge.}{$k$ pair $\rightarrow$ égalité de vote possible.}
\cor{5}{P-ALGO-04}{$V$ décroît de $6$ à $3$ $\rightarrow$ terminaison.}{cible absente.}
\cor{6}{P-ALGO-05}{$28 = 10+10+5+2+1$ (5 pièces).}{pièce 1 absente $\rightarrow$ échec possible.}
\cor{7}{P-ALGO-03}{$k=3$ $\rightarrow$ rouge~2, bleu~1 $\rightarrow$ classe rouge.}{$k$ pair $\rightarrow$ égalité.}
\cor{8}{P-ALGO-04}{sur \texttt{cible=23}, $V$ décroît $6\rightarrow 3\rightarrow 1$ $\rightarrow$ terminaison prouvée.}{\texttt{cible=38} $\rightarrow$ $V$ atteint $0$, arrêt sans trouver.}
\cor{9}{P-ALGO-03}{3 plus proches $= A(3,2), B(5,4), A(2,3)$ ; vote $A=2, B=1$ $\rightarrow$ classe $A$.}{$k=2$ avec égalité $\rightarrow$ résultat indéterminé.}

\noindent Détail des distances au nouveau point $(4,3)$ :
\begin{center}
\begin{tabular}{cccc}
\toprule
Point & Classe & Distance à $(4,3)$ & Rang \\
\midrule
$(5,4)$ & B & $\sqrt{2} \approx 1{,}41$ & 1 \\
$(3,2)$ & A & $\sqrt{2} \approx 1{,}41$ & 2 \\
$(2,3)$ & A & $\sqrt{4} = 2{,}00$ & 3 \\
$(1,1)$ & A & $\sqrt{13} \approx 3{,}61$ & 4 \\
$(8,7)$ & B & $\sqrt{32} \approx 5{,}66$ & 5 \\
\bottomrule
\end{tabular}
\end{center}
Les $k=3$ plus proches sont $B, A, A$ $\rightarrow$ vote $A=2$, $B=1$ $\rightarrow$ classe prédite $A$.

\section*{Corrigé du TP}
\begin{itemize}[nosep]
  \item Dichotomie : \texttt{cible=37} $\rightarrow$ indice $4$.
  \item Glouton : \texttt{montant=28} $\rightarrow$ $10+10+5+2+1$ (5 pièces).
  \item k-NN : \texttt{k=3} $\rightarrow$ classe rouge.
\end{itemize}

\section*{Corrigé de l'évaluation}
\begin{enumerate}[nosep,label=Q\arabic*.]
  \item milieux $18$ puis $37$ $\rightarrow$ indice $4$.
  \item $V = \text{droite} - \text{gauche} + 1$ décroît de $6$ à $3$ $\rightarrow$ terminaison.
  \item $28 = 10+10+5+2+1$ (5 pièces).
  \item rouge (2 voix) vs bleu (1) $\rightarrow$ classe rouge.
\end{enumerate}
\end{document}
